\documentclass[11pt]{article}
\usepackage[margin=1in]{geometry}
\usepackage{amsmath,amssymb}
\usepackage{xurl}
\usepackage[hidelinks]{hyperref}
\sloppy
\begin{document}
\noindent\textbf{Author:} Rick\\
\textbf{Date:} 2026-10-28\\
\textbf{To:} Clio Vega (cc Robin Langer)\\
\textbf{Title:} FPSAC 2027 cut: LG/OG negative-degree Schubert integrals --- ready for review\\
\textbf{Work-in-progress commit:} \texttt{6dbecb8}\\
\textbf{PDF at that commit:} \url{https://github.com/grandpa-rick/work-in-progress/blob/6dbecb8/papers/lg-og-fpsac-cut/main.pdf}\\
\textbf{Source:} \url{https://github.com/grandpa-rick/work-in-progress/blob/6dbecb8/papers/lg-og-fpsac-cut/main.tex}\\
\textbf{Status:} NOT submitted. Use 6dbecb8, not the earlier bac1058.

\medskip\hrule\medskip

\section*{Time-critical}
The FPSAC 2027 deadline is \textbf{2026-11-15}. Nothing is submitted without Robin's yes. Robin owns the venue yes/no and the six blue \verb|\robin{}| fields:
authorship/presenter, affiliation, human acknowledgements, AI-declaration wording, the model list, and the human role.

\section*{What the cut contains}
\begin{itemize}
\item 12 counted pages including the bibliography; the AI declaration follows on two uncounted pages. 0 TODOs.
\item The LG main line uses only (G1) and (LS): the Demazure character formula and Littelmann's path model. This covers Thm 1.1 (strict reciprocity), the recursion $(\star)$ (Prop 4.3, proved by peel + polynomiality), the integer window, the half-integer law (Thm 1.5) and Cor 1.6.
\item Buch--Ravikumar (BR) is cited in exactly two remarks. Rem 4.4 gives the Pieri values $c_\rho=(-1)^{|\theta|-1}V(\theta)$ as the $p=1$ case of their Cor 5.6, with no novelty claim. Rem 4.5 gives the OG K-theoretic sketch via their Lemma 4.5. The OG theorem itself is cited as classical.
\item Check run: a fresh script re-derives every Table 1 value and every displayed formula from the printed statements. It found 0 failures (\texttt{scripts/prove260/pdf\_check.py}).
\item Cold referee read (PROVE 260b): 0 mathematical errors. It fixed three precision points. The LS bijection now defines $a_s=1$. The $\wedge U$ Euler-characteristic identity needs (G1) for $\chi$ at all weights, not only dominant ones (the paper doesn't use it). The $(3,2,1)$ value at $-7/2$ is now described as attaining the Thm 1.5(iii) bound, because it is not the equality case.
\end{itemize}

\section*{Requests for Clio}
\begin{enumerate}
\item Line-read \S3 (LS proof) and the proof of Prop 4.3. Nobody but me has read these.
\item Line-read the compressed proof of Thm 1.5 (HZ) in \S6. It is the most compressed part.
\item Check that the AI declaration describes your reviews accurately: Thm Z (uid451), the window for both types (uid457), the $m\le 6$ census, and the 1243-shape Pieri check.
\end{enumerate}

\noindent The cut is attached as \texttt{main.pdf}, together with this note.
\end{document}
